\section{Síntesis y ampliación}

\subsection{Conclusiones principales}

\begin{importantbox}{Once conclusiones transferibles}
\begin{enumerate}
  \item El problema de investigación determina el espacio de estados, la señal de retroalimentación y la forma de la infraestructura.
  \item La intuición geométrica puede poner en marcha el aprendizaje no supervisado, pero necesita un ciclo de entrenamiento capaz de corregir las hipótesis.
  \item La corrección de una formal proof proviene del verifier; el LLM se encarga de proponer candidatos y de guiar la búsqueda de forma heurística.
  \item Decir «Optimal» exige especificar al mismo tiempo el objective, el constraint y el workload.
  \item Que algo funcione correctamente a pequeña escala no garantiza estabilidad numérica a gran escala; el entrenamiento necesita canary por etapas.
  \item El final run cost no equivale al R\&D cost; los experimentos fallidos y la construcción del sistema forman la mayor parte del iceberg.
  \item El trabajo con datos es poco visible, pero con frecuencia determina la calidad del modelo más que el código de entrenamiento.
  \item Un open-weight checkpoint es solo la primera capa de la pila de producto, no una solution utilizable.
  \item La personalización empresarial debe conectar datos, evaluación, despliegue, monitoreo y reversión.
  \item La unidad escalable del reasoning es el ciclo environment--verifier, no un texto más largo.
  \item El post-training y los activos abiertos desplazan la competencia hacia la iteración rápida, las herramientas, la retroalimentación y el control operativo.
\end{enumerate}
\end{importantbox}

\subsection{Tarea práctica: diseñar un sistema de modelo empresarial con restricciones}

Elige un escenario, por ejemplo un asistente de documentación hospitalaria, un copilot de código interno, la revisión de reclamaciones de seguros o un edge agent sin conexión, y completa el siguiente diseño:

\begin{enumerate}
  \item Escribe las restricciones no negociables: residencia de datos, latencia, disponibilidad, hardware, responsabilidad y ventanas de actualización;
  \item Elige entre external API, private cloud, on-prem o edge, y explica por qué no elegiste las demás modalidades;
  \item Define los componentes de las cinco capas de checkpoint-to-solution y su owner;
  \item Diseña una evaluation suite mínima que distinga entre errores del modelo, errores de herramientas y errores del sistema;
  \item Diseña el feedback/data loop: indica qué datos pueden reincorporarse, durante cuánto tiempo se conservan y cómo se retiran;
  \item Propón las estrategias de canary, rollback y auditoría para la actualización del base model y para las actualizaciones de post-training.
\end{enumerate}

\begin{knowledgebox}{Criterios de aceptación}
Una propuesta de alta calidad no se limita a escribir «usar Mistral / LLaMA y hacer ajuste fino». Debe poder responder: qué entorno produce la ground truth, quién verifica los resultados, cómo se obtienen los datos, cómo se sirve el modelo, cómo se localizan las fallas, cuándo se detiene el rollout y quién responde por el resultado final de negocio.
\end{knowledgebox}

\subsection{Lecturas adicionales}

\begin{enumerate}
  \item Guillaume Lample et al., \emph{Unsupervised Machine Translation Using Monolingual Corpora Only}, 2018, véase el \href{run:source-materials/unsupervised-machine-translation.pdf}{PDF local}.
  \item Guillaume Lample et al., \emph{HyperTree Proof Search for Neural Theorem Proving}, 2022, véase el \href{run:source-materials/hypertree-proof-search.pdf}{PDF local}.
  \item Jordan Hoffmann et al., \emph{Training Compute-Optimal Large Language Models}, 2022, véase el \href{run:source-materials/chinchilla-scaling-laws.pdf}{PDF local}.
  \item Hugo Touvron et al., \emph{LLaMA: Open and Efficient Foundation Language Models}, 2023, véase el \href{run:source-materials/llama.pdf}{PDF local}.
  \item Albert Q. Jiang et al., \emph{Mistral 7B}, 2023, véase el \href{run:source-materials/mistral-7b.pdf}{PDF local}.
  \item Página oficial de lanzamiento de Mistral 7B de Mistral AI: \href{https://mistral.ai/news/announcing-mistral-7b}{Mistral 7B}.
  \item Directrices y calendario de la Comisión Europea sobre GPAI: \href{https://digital-strategy.ec.europa.eu/en/policies/guidelines-gpai-providers}{Guidelines for providers of general-purpose AI models}.
\end{enumerate}

\subsection{Ruta de lectura sugerida}

Si el objetivo es reproducir la cadena de razonamiento de esta clase, y no leer cada artículo de forma aislada, se puede seguir el orden siguiente. Primero, lee el problem setup, el objetivo de entrenamiento y la ablation del artículo sobre traducción automática no supervisada, y observa cómo una hipótesis geométrica se corrige mediante denoising y back-translation; después, lee la interfaz del entorno, el hyper-tree backup y el online training de HTPS para entender cómo el verifier modifica el algoritmo de búsqueda. A continuación, lee Chinchilla y LLaMA en conjunto, escribe la restricción de presupuesto que optimiza cada uno y compara qué implica, para el despliegue de inferencia, entrenar un modelo pequeño con un horizon largo. Por último, lee las secciones de GQA/SWA de Mistral 7B y convierte la architecture claim en un balance de recursos: KV cache, ancho de banda, latencia y throughput.

\begin{knowledgebox}{Plantilla de registro de lectura}
Para cada material, registra al menos cinco elementos: \textbf{restricciones del problema}, \textbf{señal de retroalimentación autorizada}, \textbf{mecanismo central}, \textbf{principales ablaciones o evidencias} y \textbf{conclusiones que no se pueden derivar}. Al terminar, agrega una columna más: «si se llevara a production, qué componentes del sistema faltarían». Esta plantilla convierte la lectura de artículos, de un simple extracto de conclusiones, en un ejercicio de diseño de sistemas.
\end{knowledgebox}

\end{document}
